\documentclass[12pt, a4paper]{book}

\usepackage[T1]{fontenc}
\usepackage[utf8]{inputenc}
\usepackage[charter, cal=cmcal]{mathdesign}
\usepackage[scaled=0.92]{helvet}
\usepackage{geometry}
\geometry{a4paper, top=30mm, bottom=30mm, inner=32mm, outer=25mm, marginparwidth=18mm}

\usepackage{amsmath, amsthm}
\usepackage[version=4]{mhchem}
\usepackage{graphicx}
\usepackage{xcolor}
\usepackage{tikz}
\usetikzlibrary{arrows.meta, shapes.geometric, positioning, calc, backgrounds}
\usepackage{pgfplots}
\pgfplotsset{compat=1.18}
\usepackage{tcolorbox}
\tcbuselibrary{skins, breakable, theorems}
\usepackage{titlesec}
\usepackage{fancyhdr}
\usepackage{enumitem}
\usepackage{booktabs}
\usepackage{array}
\usepackage{parskip}
\usepackage{siunitx}
\usepackage[colorlinks=true, linkcolor=BrickRed,
            urlcolor=NavyBlue, citecolor=ForestGreen]{hyperref}

% ── Colours (Journey 7 — forest green / amber) ────────────────────────────────
\definecolor{JourneyColor}{HTML}{065F46}
\definecolor{ModuleColor}{HTML}{059669}
\definecolor{PuzzleColor}{HTML}{B45309}
\definecolor{DiscoveryColor}{HTML}{D97706}
\definecolor{WarningColor}{HTML}{991B1B}
\definecolor{ProblemColor}{HTML}{6B21A8}
\definecolor{ReflectColor}{HTML}{0D9488}
\definecolor{HistoryColor}{HTML}{78350F}

% ── Boxes ─────────────────────────────────────────────────────────────────────
\newtcolorbox{puzzlebox}[1][]{enhanced,breakable,
  colback=PuzzleColor!6,colframe=PuzzleColor,fonttitle=\bfseries\sffamily,
  title={The Puzzle},attach boxed title to top left={yshift=-2mm,xshift=6mm},
  boxed title style={colback=PuzzleColor},borderline west={3pt}{0pt}{PuzzleColor},#1}

\newtcolorbox{discoverybox}[1][]{enhanced,breakable,
  colback=DiscoveryColor!8,colframe=DiscoveryColor,fonttitle=\bfseries\sffamily,
  title={Discovery},attach boxed title to top left={yshift=-2mm,xshift=6mm},
  boxed title style={colback=DiscoveryColor},borderline west={3pt}{0pt}{DiscoveryColor},#1}

\newtcolorbox{problembox}[1][]{enhanced,breakable,
  colback=ProblemColor!5,colframe=ProblemColor,fonttitle=\bfseries\sffamily,
  title={Problems},attach boxed title to top left={yshift=-2mm,xshift=6mm},
  boxed title style={colback=ProblemColor},borderline west={3pt}{0pt}{ProblemColor},#1}

\newtcolorbox{deepbox}[1][]{enhanced,breakable,
  colback=JourneyColor!5,colframe=JourneyColor,fonttitle=\bfseries\sffamily,
  title={Deep Problems},attach boxed title to top left={yshift=-2mm,xshift=6mm},
  boxed title style={colback=JourneyColor},borderline west={4pt}{0pt}{JourneyColor},#1}

\newtcolorbox{historybox}[1][]{enhanced,breakable,
  colback=HistoryColor!7,colframe=HistoryColor!60,fonttitle=\bfseries\sffamily\small,
  title={Historical Note},attach boxed title to top left={yshift=-2mm,xshift=6mm},
  boxed title style={colback=HistoryColor!80},#1}

\newtcolorbox{reflectionbox}[1][]{enhanced,breakable,
  colback=ReflectColor!6,colframe=ReflectColor,fonttitle=\bfseries\sffamily,
  title={Reflection},attach boxed title to top left={yshift=-2mm,xshift=6mm},
  boxed title style={colback=ReflectColor},borderline west={3pt}{0pt}{ReflectColor},#1}

\newtcolorbox{namedconcept}[2][]{enhanced,
  colback=ModuleColor!6,colframe=ModuleColor,fonttitle=\bfseries\sffamily\large,
  title={#2},attach boxed title to top left={yshift=-2mm,xshift=6mm},
  boxed title style={colback=ModuleColor},#1}

% ── Section formatting ────────────────────────────────────────────────────────
\titleformat{\chapter}[display]
  {\normalfont\sffamily\huge\bfseries\color{JourneyColor}}
  {\textcolor{JourneyColor!60}{\chaptertitlename\ \thechapter}}{10pt}{\Huge}
  [\vspace{2pt}\textcolor{JourneyColor}{\titlerule[1.5pt]}]
\titleformat{\section}
  {\normalfont\sffamily\Large\bfseries\color{ModuleColor}}
  {\thesection}{1em}{}[\textcolor{ModuleColor!40}{\titlerule[0.6pt]}]
\titleformat{\subsection}
  {\normalfont\sffamily\large\bfseries\color{JourneyColor!80}}{\thesubsection}{1em}{}
\titleformat{\subsubsection}
  {\normalfont\sffamily\normalsize\bfseries}{\thesubsubsection}{1em}{}

\pagestyle{fancy}
\fancyhf{}
\fancyhead[LE]{\small\sffamily\textcolor{JourneyColor}{\leftmark}}
\fancyhead[RO]{\small\sffamily\textcolor{ModuleColor}{\rightmark}}
\fancyfoot[C]{\small\thepage}
\renewcommand{\headrulewidth}{0.4pt}

\newcommand{\newterm}[1]{\textit{\textcolor{DiscoveryColor}{#1}}}
\newcommand{\important}[1]{\textcolor{ModuleColor}{\textbf{#1}}}

\begin{document}

\setcounter{chapter}{6}
\chapter[Journey 7 --- Organic Chemistry: The World It Builds]
        {Journey Seven\\[8pt]
         \textcolor{ModuleColor}{\large Organic Chemistry: The World It Builds}}
\setcounter{section}{1}

\section{Module 7.2 --- Nucleic Acids and DNA Structure}
\label{sec:mod7.2}

\subsection*{Motivation}

We explored in Module 7.1 how chiral geometry and binding free energies govern molecular recognition at the clinical scale.
Now, we scale up to the absolute genetic blueprint of life: DNA.

Deoxyribonucleic acid (DNA) is a massive organic polymer.
Its primary role is informational: it stores the hereditary instructions for all living organisms.
Yet, this information is not written in ink. It is encoded in the precise 3D spatial alignment of hydrogen bonds.
This module explores the organic chemistry of nucleotides, the geometric matching rules of base pairing, and the thermodynamic forces—from hydrogen bonding to hydrophobic $\pi-\pi$ stacking—that stabilize the double helix against the thermal fluctuations of the cell.

\subsection{Watson and Crick's Cardboard Bases}

We analyze how the geometry of hydrogen bonding dictates the specificity of genetic coding.

\begin{puzzlebox}
\textbf{The base Pairing Rules.}

In 1953, James Watson and Francis Crick cut out cardboard shapes of the four nitrogenous bases: Adenine (A), Thymine (T), Guanine (G), and Cytosine (C). They tried to arrange them in pairs to fit inside a uniform double helix of width $\sim \SI{20}{\text{\AA}}$.
\begin{enumerate}[label=(\alph*), leftmargin=*]
  \item Adenine and Guanine are double-ringed structures called **purines**. Thymine and Cytosine are single-ringed structures called **pyrimidines**.
        Why would a purine-purine pair or a pyrimidine-pyrimidine pair violate the uniform \SI{20}{\text{\AA}} width constraint of the helix?
  \item Watson and Crick discovered that Adenine forms hydrogen bonds specifically with Thymine, and Guanine with Cytosine.
        * **A-T pair:** Forms two hydrogen bonds: $\ce{N-H\cdots O}$ and $\ce{N\cdots H-N}$.
        * **G-C pair:** Forms three hydrogen bonds: $\ce{O\cdots H-N}$, $\ce{N-H\cdots N}$, and $\ce{N-H\cdots O}$.
        Explain why mismatched pairings (like A-C or G-T) are geometrically unstable, failing to align the hydrogen-bond donors (electronegative atoms bonded to H) and acceptors (lone pairs on electronegative atoms).
  \item Erwin Chargaff discovered in 1950 that in all natural DNA samples, the amount of Adenine equals Thymine ($A=T$), and Guanine equals Cytosine ($G=C$). How does base-pairing geometry explain **Chargaff's rules**?
\end{enumerate}
\end{puzzlebox}

\subsection{The Phosphodiester Backbone and Directionality}

We study the polymer linkages that form the structural rails of the DNA ladder.

\begin{explobox}
Each nucleotide consists of a nitrogenous base, a deoxyribose sugar, and a phosphate group.
These nucleotides polymerize via a condensation reaction to form a **phosphodiester linkage** between the 3' (three-prime) hydroxyl group of one sugar and the 5' (five-prime) phosphate group of the next.
\begin{enumerate}[label=(\alph*), leftmargin=*]
  \item Explain why the phosphodiester backbone gives DNA a clear chemical directionality (meaning a strand runs from the 5' end to the 3' end).
  \item The two strands in a double helix run in opposite directions; we call them **antiparallel**.
        If one strand has the sequence 5'-ATCG-3', write the sequence of the complementary strand, indicating its 5' and 3' directionality.
  \item The phosphate groups in the backbone are completely ionized at physiological pH, carrying a negative charge.
        Since like charges repel, why doesn't the double helix immediately fly apart? (Think about the role of dissolved positive ions like $\ce{Na+}$ and $\ce{Mg^2+}$ in cell water).
\end{enumerate}
\end{explobox}

\subsection{Helix Stabilization: The Aromatic Stack}

We explore the thermodynamic forces that keep the double helix bound.

\begin{explobox}
A common misconception is that the hydrogen bonds between bases are the primary force holding the two strands together.
In water, water molecules form strong hydrogen bonds with the bases when the DNA is single-stranded. Thus, the net energy change ($\Delta H$) of forming base-pair hydrogen bonds compared to base-water hydrogen bonds is close to zero.
\problem The primary thermodynamic force stabilizing the double helix is actually **hydrophobic $\pi-\pi$ stacking**.
\begin{enumerate}[label=(\alph*), leftmargin=*]
  \item The nitrogenous bases are highly planar, aromatic rings (from Journey 6 Module 3).
        Are these aromatic rings hydrophilic or hydrophobic?
  \item In the double helix, the bases are stacked flat on top of each other, like a stack of coins.
        Explain how this stacking minimizes the contact of these hydrophobic surfaces with water, releasing water molecules into the bulk solution to increase entropy ($\Delta S > 0$).
  \item How do the delocalized $\pi$ electron clouds of the stacked bases interact with each other to provide dispersion forces (Van der Waals interactions)?
\end{enumerate}
\end{explobox}

\newpage

% ─── Section: Discovery ──────────────────────────────────────────────────────
\section{Part 3 --- Discovery}

\begin{discoverybox}
\textbf{What have we built?}

\textbf{First.} The specificity of genetic encoding arises from the geometry of hydrogen bonding: Adenine pairs with Thymine via two hydrogen bonds, and Guanine pairs with Cytosine via three hydrogen bonds.

\textbf{Second.} The phosphodiester backbone links nucleotides from the 5' phosphate to the 3' hydroxyl, forming two complementary, antiparallel strands.

\textbf{Third.} While hydrogen bonds provide the specificity required for exact replication, the primary thermodynamic stabilization of the double helix comes from hydrophobic $\pi-\pi$ stacking interactions between stacked aromatic bases.

\textbf{Fourth.} The electrostatic repulsion between negatively charged phosphate backbones is screened in vivo by counter-ions like magnesium.
\end{discoverybox}

\newpage

% ─── Section: Naming Concepts ────────────────────────────────────────────────
\section{Part 4 --- What This Is Called}

\begin{namedconcept}{Nucleic acid parameters}
* **Nucleotide:** The monomeric building block of nucleic acids, consisting of a nitrogenous base, a five-carbon sugar, and a phosphate group.
* **Chargaff's Rules:** The observation that in double-stranded DNA, the ratio of A to T is 1:1, and the ratio of G to C is 1:1.
* **Antiparallel:** The relative orientation of the two strands in a double helix, running in opposite chemical directions (5' $\to$ 3' and 3' $\to$ 5').
* **$\pi-\pi$ Stacking:** The attractive non-covalent interaction between planar aromatic rings, driven by hydrophobic effects and London dispersion forces.
* **Melting Temperature ($T_m$):** The temperature at which exactly half of a double-stranded DNA sample is denatured (melted) into single strands.
\end{namedconcept}

\newpage

% ─── Section: Solved Exercises ───────────────────────────────────────────────
\section{Part 5 --- Solved Exercises}

\begin{problembox}
\begin{enumerate}[label=\textbf{E\arabic*.}]
  \item \textbf{Applying Chargaff's Rules.}
  
        A double-stranded DNA sample isolated from an extreme thermophilic bacterium contains $35.0\%$ Guanine (G). Calculate the percentages of Adenine (A), Thymine (T), and Cytosine (C) in this sample.
        
        \textbf{Solution:}
        By Chargaff's rules, the percentage of Cytosine must equal Guanine:
        \[
          \% \ce{C} \;=\; \% \ce{G} \;=\; 35.0\%.
        \]
        The sum of all four bases must equal $100.0\%$:
        \[
          \% \ce{A} + \% \ce{T} + \% \ce{G} + \% \ce{C} \;=\; 100.0\%.
        \]
        Substitute the values for G and C:
        \[
          \% \ce{A} + \% \ce{T} + 35.0\% + 35.0\% \;=\; 100.0\% \implies \% \ce{A} + \% \ce{T} + 70.0\% \;=\; 100.0\%.
        \]
        This leaves $30.0\%$ for Adenine and Thymine:
        \[
          \% \ce{A} + \% \ce{T} \;=\; 30.0\%.
        \]
        By Chargaff's rules, the percentage of Adenine must equal Thymine:
        \[
          \% \ce{A} \;=\; \% \ce{T} \implies 2 (\% \ce{A}) \;=\; 30.0\% \implies \% \ce{A} \;=\; 15.0\%.
        \]
        Thus, the total percentages are:
        * **Guanine:** $35.0\%$
        * **Cytosine:** $35.0\%$
        * **Adenine:** $15.0\%$
        * **Thymine:** $15.0\%$
        
  \item \textbf{Length of the Human Genome.}
  
        In the B-form double helix, the rise per base pair (the vertical distance between stacked base pairs along the helix axis) is $d \approx \SI{3.4}{\text{\AA}}$ ($3.4 \times 10^{-10}$ m).
        The haploid human genome contains $N = 3.2 \times 10^9$ base pairs. Calculate the physical length of a single haploid genome's DNA in meters.
        
        \textbf{Solution:}
        The total physical length $L$ of the DNA is the product of the number of base pairs and the rise per base pair:
        \[
          L \;=\; N \cdot d.
        \]
        Substitute the given values:
        \[
          L \;=\; \left( 3.2 \times 10^9 \text{ base pairs} \right) \left( 3.4 \times 10^{-10} \text{ m/base pair} \right)
        \]
        \[
          L \;=\; 3.2 \times 3.4 \times 10^{-1} \text{ m} \;=\; 10.88 \times 10^{-1} \text{ m} \;=\; 1.09 \text{ m}.
        \]
        Thus, a single haploid genome contains about \SI{1.09}{m} of DNA. If we consider a diploid human cell (containing two haploid copies), there is over \SI{2}{m} of DNA packed inside a cell nucleus only a few micrometers wide!
\end{enumerate}
\end{problembox}

\newpage

% ─── Section: Problems ───────────────────────────────────────────────────────
\section{Part 6 --- Problems}

\begin{problembox}
\begin{enumerate}[label=\textbf{P\arabic*.}]
  \item \textbf{Geometries of Base Pairing.}
  
        Consider the chemical structures of the A-T and G-C base pairs.
        \begin{enumerate}[label=(\alph*)]
          \item Draw a diagram showing the hydrogen bonds in a G-C base pair. Label each hydrogen bond, indicating the electronegative donor atom (with its covalent H) and the acceptor atom (with its lone pair).
          \item When we heat DNA in water, the thermal motion breaks the interactions, separating the double helix into single strands (denaturation).
                Explain why DNA containing a higher percentage of G-C base pairs requires a higher temperature to melt (denature) than DNA with a lower G-C content.
        \end{enumerate}

  \item \textbf{Single-Stranded Genomes.}
  
        The DNA genome of the bacterial virus (bacteriophage) $\phi$X174 was isolated and analyzed, yielding the following base composition:
        \[
          \% \ce{A} \;=\; 24.0\%, \qquad \% \ce{T} \;=\; 33.0\%, \qquad \% \ce{G} \;=\; 23.0\%, \qquad \% \ce{C} \;=\; 20.0\%.
        \]
        \begin{enumerate}[label=(\alph*)]
          \item Show why this composition violates Chargaff's rules.
          \item What does this violation suggest about the physical structure of the DNA inside this virus?
          \item If this virus infects a bacterium, the host machinery copies the viral DNA to form a double-stranded replicative form. Calculate the base composition of this double-stranded intermediate.
        \end{enumerate}
\end{enumerate}
\end{problembox}

\newpage

% ─── Section: Deep Problems ──────────────────────────────────────────────────
\section{Part 7 --- Deep Problems}

\begin{deepbox}
\begin{enumerate}[label=\textbf{D\arabic*.}]
  \item \textbf{Thermodynamics of DNA Melting.}
  
        The transition between double-stranded DNA ($\ce{dsDNA}$) and single-stranded DNA ($\ce{ssDNA}$) is a dynamic equilibrium:
        \[
          \ce{dsDNA <=> 2 ssDNA}.
        \]
        Let $\theta$ be the fraction of denatured DNA ($\theta = 0$ for fully double-stranded, $\theta = 1$ for fully single-stranded).
        The melting temperature $T_m$ is the temperature at which $\theta = 0.5$.
        \begin{enumerate}[label=(\alph*)]
          \item The equilibrium constant at $T_m$ is related to the standard enthalpy ($\Delta H^\circ$) and entropy ($\Delta S^\circ$) of denaturation by $\Delta G^\circ = \Delta H^\circ - T_m \Delta S^\circ = 0$.
                Prove that:
                \[
                  T_m \;=\; \frac{\Delta H^\circ}{\Delta S^\circ}.
                \]
          \item Plot the fraction of denatured DNA $\theta$ versus temperature $T$. Why is this denaturation curve highly sigmoidal (a cooperative transition, where breaking the first few base pairs makes it much easier to break the rest)?
          \item How does the difference in hydrogen bonding (3 bonds in G-C vs. 2 in A-T) affect the denaturation enthalpy $\Delta H^\circ$, and how does this mathematically shift the melting temperature $T_m$ according to your formula?
        \end{enumerate}

  \item \textbf{Intercalating Mutagens and $\pi-\pi$ Stacking.}
  
        Some planar, aromatic molecules, such as ethidium bromide ($\ce{EtBr}$) and proflavine, are potent mutagens.
        They act by sliding directly in between adjacent base pairs in the double helix, a process called **intercalation**.
        \begin{enumerate}[label=(\alph*)]
          \item Explain why the planar, aromatic structure of these molecules allows them to insert into the hydrophobic core of the double helix. What forces drive this insertion?
          \item When an intercalator inserts, it increases the distance between adjacent base pairs from \SI{3.4}{\text{\AA}} to \SI{6.8}{\text{\AA}}, unwinding and deforming the helix geometry.
                Explain how this physical deformation leads to insertion or deletion mutations (frameshifts) when the DNA polymerase enzyme tries to copy the deformed template.
        \end{enumerate}
\end{enumerate}
\end{deepbox}

\newpage

% ─── Section: Reflection ─────────────────────────────────────────────────────
\section{Part 8 --- Reflection}

\begin{reflectionbox}
\textbf{Conceptual Reflection}

Consider the following questions:
\begin{enumerate}
  \item **On the Stability of Information:** If DNA was stabilized by covalent bonds between the two strands instead of weak non-covalent forces (hydrogen bonds and $\pi-\pi$ stacking), how would the cell copy the information? Why is the weak, reversible nature of these non-covalent bonds essential for genetic replication and transcription?
  \item **On Molecular Packaging:** Human cells pack \SI{2}{m} of DNA into a nucleus only \SI{5}{\micro m} wide. If the DNA backbone was achiral and lacked the helical twist, would it be possible to coil and compress the molecule so densely without tangling?
  \item **On Biophysics:** We saw that the sigmoidal shape of the DNA melting curve is a cooperative transition—meaning it behaves like a phase transition (like water boiling). How does the cooperative behavior of $10^6$ weak chemical bonds acting together create a sharp, binary threshold in biological state?
\end{enumerate}
\end{reflectionbox}

\medskip
\noindent\textbf{Looking ahead.}
Module~7.3 completes this journey by introducing **Proteins and Molecular Machines**. Having analyzed the DNA blueprint, we will explore the active engines that DNA codes for. We will study peptide bonds, protein folding, and examine the rotary engine of ATP Synthase.

\vspace{1.4cm}
\begin{center}
{\color{JourneyColor}\rule{0.5\linewidth}{0.8pt}}\\[6pt]
{\small\color{JourneyColor}\textit{End of Journey 7, Module 7.2 --- Nucleic Acids and DNA Structure}}\\[4pt]
{\small\color{JourneyColor}Continue to: \textbf{Module 7.3 --- Proteins and Molecular Machines}}\\[6pt]
{\color{JourneyColor}\rule{0.5\linewidth}{0.8pt}}
\end{center}

\end{document}
